% BHU PYQ -- Transcribed & Corrected
% Paper: CSCMJ42: Data Base Management Systems
% Exam: B.Sc. (Hons.) Semester IV Examination, 2025-26 (NEP)
% Category: Major Course
% Department: Computer Science

\documentclass[11pt,a4paper]{article}
\usepackage[a4paper,top=2.5cm,bottom=2.5cm,left=2.8cm,right=2.8cm,headheight=14pt]{geometry}
\usepackage{amsmath,amssymb}
\usepackage[T1]{fontenc}
\usepackage[utf8]{inputenc}
\usepackage{lmodern,microtype,fancyhdr,enumitem,hyperref,lastpage,parskip}

\hypersetup{
    colorlinks=true,
    urlcolor=black,
    linkcolor=black,
    pdftitle={CSCMJ42: Data Base Management Systems -- BHU Sem IV 2025-26 (NEP)}
}

\pagestyle{fancy}
\fancyhf{}
\renewcommand{\headrulewidth}{0.4pt}
\renewcommand{\footrulewidth}{0.4pt}
\fancyhead[L]{\small   \textit{Banaras Hindu University}}
\fancyhead[R]{\small   \textit{CSCMJ42: Database Management Systems (Sem IV 2025-26)}}
\fancyfoot[C]{\small Page  \thepage\ of \pageref{LastPage}}


\newlist{parts}{enumerate}{2}
\setlist[parts,1]{label=(\alph*),leftmargin=*,itemsep=4pt,topsep=2pt}
\setlist[parts,2]{label=(\roman*),leftmargin=*,itemsep=2pt,topsep=2pt}

\newcommand{\pts}[1]{\hfill   \textit{[#1]}}


\begin{document}


\begin{center}
  {\large\bfseries BANARAS HINDU UNIVERSITY}\\[2pt]
  {\normalsize B.Sc. (Hons.) --- Semester IV Examination, 2025-26 (NEP)}\\[6pt]
  \rule{0.8\linewidth}{0.5pt}\\[6pt]
  {\large\bfseries COMPUTER SCIENCE}\\[4pt]
  {\large\bfseries Paper Code: CSCMJ42 : Data Base Management Systems}\\[4pt]
  {\normalsize Time: 2:30 Hours\hfill Maximum Marks: 70}
\end{center}

\rule{\linewidth}{0.4pt}

\smallskip



\noindent   \textit{Instructions: Answer any   \textbf{five} questions.}

\bigskip




\noindent   \textbf{Question 1.}

\begin{parts}
  \item Write short notes on any three of the following:\pts{6}
  
\begin{parts}
    \item Data Warehouse
    \item Functionalities and characteristics of DBMS
    \item Database users
    \item Recursive relationship type in ER model
  \end{parts}
  \item What do you understand by data models? Write different categories of data models.\pts{4}
  \item Explain the following:\pts{4}
  
\begin{parts}
    \item Key constraint
    \item Entity integrity constraints
    \item Referential integrity constraint
    \item Disjointness and completeness constraint on specialization/generalization
  \end{parts}
\end{parts}

\medskip\rule{\linewidth}{0.2pt}\medskip




\noindent   \textbf{Question 2.}

\begin{parts}
  \item Give examples of the following types of specialization/generalization:\pts{8}
  
\begin{parts}
    \item Disjoint, Total
    \item Disjoint, Partial
    \item Overlapping, Total
    \item Overlapping, Partial
  \end{parts}
  \item Explain Theta JOIN, EQUIJOIN and NATURAL JOIN with suitable examples.\pts{6}
\end{parts}

\medskip\rule{\linewidth}{0.2pt}\medskip




\noindent   \textbf{Question 3.}

\begin{parts}
  \item Explain unary relational algebra operations with suitable examples of each.\pts{9}
  \item Write relational algebra operations for the following queries:\pts{5}
  
\begin{parts}
    \item Compute the count (number) of employees and their average salary.
    \item Compute the count of employees and average salary per department.
  \end{parts}
\end{parts}

\medskip\rule{\linewidth}{0.2pt}\medskip




\noindent   \textbf{Question 4.}

\begin{parts}
  \item Explain the concept of views and indexes in SQL. How are they created and used?\pts{8}
  \item Explain different anomalies (update, insert and delete) with their examples.\pts{6}
\end{parts}

\medskip\rule{\linewidth}{0.2pt}\medskip




\noindent   \textbf{Question 5.}

\begin{parts}
  \item Define Query Optimization and explain its objectives.\pts{4}
  \item Explain the working of Cost-Based Query Optimization and Heuristic Query Optimization. Describe the factors used in cost estimation.\pts{10}
\end{parts}

\medskip\rule{\linewidth}{0.2pt}\medskip




\noindent   \textbf{Question 6.}

\begin{parts}
  \item Explain minimal sets of Functional Dependencies. What do you understand by Full functional dependency? Give an example of Full functional dependency.\pts{6}
  \item Consider the relation $R(A, B, C, D, E, F)$ with functional dependencies $F = \{A  o BC, C  o D, BD  o E, E  o F\}$. Answer the following:\pts{8}
  
\begin{parts}
    \item Find all candidate keys
    \item Normalize the relation up to BCNF.
  \end{parts}
\end{parts}

\medskip\rule{\linewidth}{0.2pt}\medskip




\noindent   \textbf{Question 7.}

\begin{parts}
  \item Explain Lost update and Temporary update problems with their examples in the context of concurrent execution of transactions. Draw state transition diagram to illustrate the states for transaction execution.\pts{7}
  \item What do you understand by conflict serializable schedule? Write an algorithm to test for conflict serializability.\pts{7}
\end{parts}

\vfill
\rule{\linewidth}{0.4pt}

\begin{center}\small
\href{https://bhu-exam.vercel.app/}{Click here to visit our website}\\[4pt]
\href{https://me-aryan.vercel.app/}{Click here to visit our contributor}
\end{center}

\end{document}
